\section{结论 (Conclusion)}
\label{sec:conclusion}

我们引入了一种新颖的计算模型，该模型能够利用现有的加密经济机制来实现可扩展性的重大改进，而不会造成持久的状态分割从而牺牲整体凝聚力。我们将此整体模式称为 collect-refine-join-accumulate。此外，我们形式化地定义了该逻辑的链上部分，本质上是 join-accumulate 部分。我们将此协议称为 \Jam 链。

我们论证 \Jam 的模型提供了一个新颖的"甜蜜点"，与完全同步模型相比，允许在安全、有弹性的共识中完成大量计算，并且与持久分割模型不同，仍然对计算的时序和集成到某个单例状态机器中具有严格的保证。

\subsection{后续工作 (Further Work)}

虽然我们能够在一些基本假设下估算理论可计算量，甚至可以与现有系统进行广泛比较，但实际数据是无价的。我们认为该模型值得进一步的实证研究，以更好地理解这些理论限制如何转化为现实世界的性能。我们认为，与现有协议进行适当的成本分析和比较也将是进一步工作的绝佳主题。

我们可以合理地确信 \Jam 的设计使其能够承载一个 Polkadot \emph{parachains} 可以在其上得到验证的服务，然而需要进一步的原型工作来了解 \textsc{pvm} 驱动的计量系统所能支持的可能吞吐量。我们将此类报告留作后续工作。同样，我们也故意省略了更高层次协议元素的细节，包括加密货币、coretime 销售、质押和常规智能合约功能。

正在考虑对本文所述协议进行若干潜在修改，以使协议的实际使用更加便利。这些包括：

\begin{itemize}
  \item accumulate 中服务之间的同步调用。
  \item 对 \texttt{transfer} 函数的限制，以允许 accumulate 过程中的大规模并行。
  \item 在某些条件下在 accumulate 期间预留大量额外计算能力的可能性。
  \item 将 Merklization 引入 Work Package 格式，以消除为评估其授权而需要下载整个包的需要。
\end{itemize}

网络协议在此阶段也被故意保持未定义状态，其描述必须在后续提案中完成。

验证者性能目前不在链上追踪。我们确实期望在 \Jam 协议的最终版本中在链上追踪，但其具体格式尚不确定，因此目前予以省略。
